% ------------------------------------------------------------------------
% USTAWIENIA
% ------------------------------------------------------------------------

% ------------------------------------------------------------------------
%   Kropki po numerach sekcji, podsekcji, itd.
% ------------------------------------------------------------------------
\makeatletter
    \def\numberline#1{\hb@xt@\@tempdima{#1.\hfil}}
    \renewcommand*\@seccntformat[1]{\csname the#1\endcsname.\enspace}
\makeatother

% ------------------------------------------------------------------------
%   Numeracja równań, rysunków i tabel (np. 1.2)
% ------------------------------------------------------------------------
\makeatletter
    \@addtoreset{equation}{section}
    \renewcommand{\theequation}{{\thesection}.\@arabic\c@equation}

    \@addtoreset{figure}{section}
    \renewcommand{\thefigure}{{\thesection}.\@arabic\c@figure}

    \@addtoreset{table}{section}
    \renewcommand{\thetable}{{\thesection}.\@arabic\c@table}
\makeatother

% ------------------------------------------------------------------------
% Środowisko tabeli
% ------------------------------------------------------------------------
\newenvironment{tabela}[3]
{
    \begin{table}[!htb]
    \centering
    \caption[#1]{#2}
    \vskip 9pt
    #3
}{
    \end{table}
}

% ------------------------------------------------------------------------
% Dostosowanie listy TODO
% ------------------------------------------------------------------------
\renewcommand{\todoitem}[2]{%
    \item \label{todo:\thetodo}%
    \ifx#1\todomark%
        \else\textbf{#1 }%
    \fi%
    (str.~\pageref{todopage:\thetodo})\ #2}
\renewcommand{\todoname}{Do zrobienia...}
\renewcommand{\todomark}{~uzupełnić}

% ------------------------------------------------------------------------
% Sekcje nienumerowane
% ------------------------------------------------------------------------
\def\nonumsection#1{%
    \section*{#1}%
    \addcontentsline{toc}{section}{#1}%
}
\def\nonumsubsection#1{%
    \subsection*{#1}%
    \addcontentsline{toc}{subsection}{#1}%
}

% ------------------------------------------------------------------------
% Inne ustawienia
% ------------------------------------------------------------------------
\frenchspacing
\setlength{\parskip}{3pt}
\linespread{1.3}
\setcounter{tocdepth}{4}
\setcounter{secnumdepth}{4}

\titleformat{\paragraph}[hang]
{\normalfont\sffamily\bfseries}{\theparagraph}{1em}{}



\sloppy


\DeclareFieldFormat{url}{\url{#1}}
